% ==== TEMPLATE LaTex DISSERTAÇÃO DE MESTRADO MEDVET 2019 =====%

% Não alterar os comandos dessa página. O template só dever ser alterado, com os devidos conhecimentos, a partir da linha 50. %

% Todas as linhas que possam a vir ser editadas estarão indicadas com comentários na frente %

\documentclass[12pt, oneside, a74paper]{article} 
\usepackage[a4paper,top=3cm,bottom=2cm,right=2cm,left=3cm]{geometry}
\usepackage[brazil]{babel}
\usepackage{amstext}
\usepackage{latexsym}
\usepackage{amsmath}
\usepackage{amssymb}
\usepackage{appendix}
\usepackage{amsfonts}
\usepackage{graphicx}
\usepackage{makeidx}
\usepackage{setspace}
\usepackage[normalem]{ulem}
\usepackage{subfigure}
\usepackage{acronym}
\usepackage{indentfirst}
\usepackage[comma, authoryear]{natbib}
\usepackage{array}
\usepackage{multirow}
\usepackage{longtable}
\usepackage[table]{xcolor}

\usepackage[export]{adjustbox}
\usepackage{wrapfig}

\usepackage{fontspec}
\setmainfont[BoldFont=calibriBold.ttf]{calibri.ttf}

\usepackage{titlesec}
\titleformat*{\section}{\fontsize{12}{12}\bfseries}
\titleformat*{\subsection}{\fontsize{12}{12}\bfseries}
\titleformat*{\subsubsection}{\fontsize{12}{12}\bfseries}
\titleformat*{\paragraph}{\fontsize{12}{12}\bfseries}
\titleformat*{\subparagraph}{\fontsize{12}{12}\bfseries}

\linespread{1.5}
\setlength{\parindent}{1.5cm}

\makeindex

\begin{document}

\input{capa.tex}%IMPORTANDO FOLHA DA CAPA

\input{folhaRosto.tex}%IMPORTANDO FOLHA DE ROSTO

\input{folhaAvaliacao.tex}%IMPORTANDO FOLHA DE AVALIAÇÃO
    
\input{dedicatoria.tex}%IMPORTANDO FOLHA DE DEDICATÓRIA. Caso não haja dedicatória, comente a linha com um símbolo de porcentagem '%' no começo.

\input{agradecimentos.tex}%IMPORTANDO FOLHA DE AGRADECIMENTOS

\input{epigrafe.tex}%IMPORTANDO FOLHA DE EPÍGRAFE. Caso não haja epígrafe, comente a linha com um símbolo de porcentagem '%' no começo.

\input{resumo.tex}%IMPORTANDO FOLHA DE RESUMO

\input{abstract.tex}%IMPORTANDO FOLHA DE ABSTRACT

%=============== FOLHA LISTA DE FIGURAS =============================%

\newgeometry{top=3.5cm, bottom=2cm, right=2cm, left=3cm}
\clearpage
\renewcommand{\listfigurename}{\begin{center} \fontsize{12}{12} \textbf{Lista de Figuras} \end{center}}
\listoffigures

%=============== FIM DA FOLHA LISTA DE FIGURAS =====================%

%=============== FOLHA LISTA DE TABELAS ============================%

\newgeometry{top=3.5cm,bottom=2cm,right=2cm,left=3cm}
\clearpage
\renewcommand{\listtablename}{\begin{center} \fontsize{12}{12} \textbf{Lista de Tabelas} \end{center}}
\listoftables

%=============== FIM DA FOLHA LISTA DE TABELAS =====================%

\input{siglas.tex} %IMPORTANDO FOLHA LISTA DE SIGLAS

\input{simbolos.tex} %IMPORTANDO FOLHA LISTA DE SÍMBOLOS

%=============== FOLHA DE SUMÁRIO =================================%
\newgeometry{top=3.5cm,bottom=2cm,right=2cm,left=3cm}
\clearpage
\renewcommand{\contentsname}{\begin{center} \fontsize{12}{12} \textbf{Sumário} \end{center}}
\tableofcontents
% =============== FIM DA FOLHA DE SUMÁRIO =========================== %

% ============== CONTEÚDO DA DISSERTAÇÃO ============================ %

% Usaremos o comando \input{} para juntarmos as partes do trabalho semelhantemente a como fizemos no artigo interno. Contudo, como nossas seções estão dentro da pasta 'sections' todos os nossos inputs começarão com sections/. Exemplo: %

% \input{sections/conclusao.tex} %

% Os comandos inputs devem ser postos na ordem de compilação do trabalho, ou seja, primeiro a introdução seguida de objetivos e etc. %

% O artigo interno tem somente um input, que é o (artigo.tex), suas subseções serão acopladas automaticamente. %

\input{TabelaFiguras.tex}% Importando section de instrução. Após a leitura da mesma, comente colocando o símbolo de porcentagem '%' no começo da linha. % 


% ============== FIM DO CONTEÚDO DA DISSERTAÇÃO ===================== %

\input{biblio.tex}

\end{document}